\chapter{Cronograma}

\section{Cronograma de ejecución de actividades}

El plan de trabajo se ha diseñado para optimizar los tiempos de intervención en los laboratorios y aulas, priorizando los ambientes críticos del Tercer Nivel (Alta Densidad).

\begin{table}[H]
\centering
\caption{Cronograma de actividades para el proyecto de cableado estructurado}
\begin{tabular}{|c|l|c|p{6.5cm}|}
\hline
\textbf{Fase} & \textbf{Actividad Principal} & \textbf{Duración Est.} & \textbf{Descripción Técnica} \\ \hline
I & Obras Preliminares y Replanteo & 3 Días & Trazo de rutas en los 4 niveles, marcado de ubicación de puntos en pared/piso y movilización de materiales a almacén en obra. \\ \hline
II & Instalación de Canalizaciones & 10 Días & Montaje de canaletas, tuberías PVC y cajas de paso desde los montantes hasta el área de usuario. Se iniciará por el Nivel 3 (Laboratorios). \\ \hline
III & Tendido de Cableado Horizontal & 8 Días & Instalación de bobinas Cat 6A U/UTP. Se realiza el peinado de mazos y etiquetado provisional en extremos. \\ \hline
IV & Conectorización y Terminación & 7 Días & Terminación de Jacks RJ45 en lado usuario y remate en Patch Panels (Gabinetes IDF/MDF). \\ \hline
V & Certificación y Etiquetado Final & 4 Días & Pruebas con equipo certificador (Fluke), colocación de etiquetas definitivas bajo norma TIA-606-C y limpieza de obra. \\ \hline
VI & Entrega de Dossier (Cierre) & 3 Días & Elaboración de planos As-Built y entrega de informe de certificación. \\ \hline
\end{tabular}
\end{table}

Nota: Algunas actividades como la Canalización (II) y el Tendido de Cable (III) se realizarán con un traslape (superposición) de tiempos para agilizar el avance.

\section{Cronograma de compras y desembolsos}

Este cronograma financiero asegura la disponibilidad de los insumos críticos (especialmente el cobre, que es sensible a variaciones de precio) antes del inicio de las labores fuertes.

\begin{table}[H]
\centering
\caption{Plan de desembolsos y cronograma de pagos del proyecto}
\begin{tabular}{|l|c|c|p{6.5cm}|}
\hline
\textbf{Hito del Proyecto} & \textbf{Momento} & \textbf{\% Desembolso} & \textbf{Concepto / Justificación} \\ \hline
Adelanto Inicial & Día 0 & 60\% & \textbf{Compra de Materiales Importados:} Adquisición del 100\% de bobinas Cat 6A, Jacks, Patch Panels y Canaletas para asegurar stock y congelar precios (mitigación de riesgo inflacionario). \\ \hline
Valorización & Día 15 & 25\% & \textbf{Avance de Obra:} Pago de planillas de personal (Mano de Obra) e insumos menores consumidos durante la canalización y tendido. \\ \hline
Liquidación & Día 30 & 15\% & \textbf{Cierre y Conformidad:} Pago final contra entrega de Protocolos de Certificación, Planos As-Built y operatividad al 100\%. \\ \hline
\end{tabular}
\end{table}

\section{Diagrama de Gantt de instalación}

A continuación se presenta la distribución temporal de las tareas durante las 4 semanas de ejecución proyectada.

\subsection*{Estrategia de intervención por niveles}

Para cumplir con este Gantt, se propone el siguiente orden de ataque:

\begin{enumerate}
    \item \textbf{Prioridad 1 (Semana 1-2): Nivel 3.} Al tener la mayor densidad (135 puntos) y ser laboratorios de cómputo, es la ruta crítica. Se debe liberar primero para no entorpecer clases.
    \item \textbf{Prioridad 2 (Semana 2-3): Nivel 2.} Nivel administrativo y biblioteca.
    \item \textbf{Prioridad 3 (Semana 3): Nivel 1 y Nivel 4.} Se ejecutan en paralelo al cierre del nivel 2, ya que tienen menor densidad.
\end{enumerate}
